\clearpage\thispagestyle{empty}
\pdfbookmark[0]{Quellen zum Bericht über Pallas und Juno vom März 1806}{pallasjuno180603-sources}
\begin{center}\textbf{Quellen zum Bericht über Pallas und Juno vom März 1806}\end{center}
\noindent Diese KI-Übertragung in heutiges Deutsch folgt dem Auszug aus Franz Xaver von Zachs Bericht in \href{https://gdz.sub.uni-goettingen.de/id/PPN236007467}{Gauß’ \textit{Werken}, Band~VI} (Göttingen, 1874), S.~269--270. Der Bericht erschien als Beitrag~XXIV, \textit{Fortgesetzte Nachrichten über die neuen Planeten}, in der \textit{Monatlichen Correspondenz}, Band~XIII, März 1806, auf \href{https://www.digitale-sammlungen.de/en/view/bsb10538605?page=323}{S.~313}, \href{https://www.digitale-sammlungen.de/en/view/bsb10538605?page=324}{S.~314} und \href{https://www.digitale-sammlungen.de/en/view/bsb10538605?page=325}{S.~315}. Die vier Beobachtungen der Pallas und die zwei der Juno auf S.~314 der Zeitschrift stehen in der Werkausgabe auf S.~267 im \hyperlink{native-272}{Brief von Gauß an Bode}. Der Bericht zitiert Carl Friedrich Gauß zu den Störungen und zu den Schwierigkeiten bei der Beobachtung der Juno. Übertragung und Satz wurden mit KI erstellt. Die Zeitschriftendigitalisate stammen von der Bayerischen Staatsbibliothek.
\clearpage
